% =====================================================================
%  PREÁMBULO — paquetes y configuración global
%  El \documentclass va en main.tex, NO aquí.
%  Stack probado en Overleaf: compilador pdfLaTeX + motor Biber.
% =====================================================================

% --- Codificación de entrada y de fuente -----------------------------
% utf8 permite escribir acentos y ñ directamente en los .tex.
% T1 evita que las palabras con acento se partan mal al final de línea.
\usepackage[utf8]{inputenc}
\usepackage[T1]{fontenc}

% --- Tipografía: Times New Roman -------------------------------------
% mathptmx sustituye la fuente por Times (requisito típico de informes
% universitarios). Si la universidad exige ARIAL, hay que cambiar el
% compilador a XeLaTeX y usar fontspec con \setmainfont{Arial}.
\usepackage{mathptmx}

% --- Márgenes --------------------------------------------------------
% 3 cm a la izquierda (para el anillado/empastado), 2,5 cm el resto.
\usepackage[a4paper,left=3cm,right=2.5cm,top=2.5cm,bottom=2.5cm]{geometry}

% --- Idioma español --------------------------------------------------
% es-tabla        -> los flotantes se rotulan «Tabla», no «Cuadro».
% es-nodecimaldot -> no fuerza la coma decimal en modo matemático.
% es-lcroman      -> numera los preliminares en romanos MINÚSCULOS (i, ii,
%                    iii), que es la convención académica. Sin esta opción,
%                    babel-spanish los pone en mayúsculas (I, II, III).
\usepackage[spanish,es-tabla,es-nodecimaldot,es-lcroman]{babel}

% --- Interlineado ----------------------------------------------------
% ÚNICO lugar donde se controla el interlineado de todo el documento.
% \onehalfspacing = 1,5 líneas (el que estamos usando).
% Si el docente pidiera interlineado doble, cambia a \doublespacing.
% Si necesitas comprimir para entrar en un límite de páginas, \singlespacing.
\usepackage{setspace}
\onehalfspacing

% --- Imágenes --------------------------------------------------------
% graphicspath evita tener que escribir la ruta completa en \includegraphics:
% basta con \includegraphics{mapa_ubicacion.png}.
\usepackage{graphicx}
\graphicspath{{imagenes/}{imagenes/mapas/}{imagenes/diagramas/}{imagenes/fotos/}{imagenes/logos/}}

% --- Tablas ----------------------------------------------------------
% booktabs  -> líneas horizontales de calidad (\toprule, \midrule, \bottomrule).
% longtable -> tablas que se parten en varias páginas (la matriz de no
%              conformidades del Anexo B seguro la necesita).
% tabularx  -> columnas de ancho automático que ajustan al ancho de página.
% multirow  -> celdas que ocupan varias filas.
% array     -> tipos de columna personalizados (ver config/comandos.tex).
\usepackage{booktabs}
\usepackage{longtable}
\usepackage{tabularx}
\usepackage{multirow}
\usepackage{array}

% --- Leyendas de tablas y figuras ------------------------------------
% Norma APA / uso académico: el rótulo de la TABLA va ARRIBA y el de la
% FIGURA va ABAJO. Aquí se fija ese comportamiento para todo el documento.
\usepackage{caption}
\captionsetup[table]{position=top,labelsep=period,font=small,labelfont=bf}
\captionsetup[figure]{position=bottom,labelsep=period,font=small,labelfont=bf}

% --- Listas ----------------------------------------------------------
% enumitem permite controlar sangrías y separación de itemize/enumerate.
\usepackage{enumitem}
\setlist{itemsep=2pt,parsep=0pt,topsep=4pt}

% --- Posicionamiento de flotantes ------------------------------------
% float habilita el especificador [H] = «aquí exactamente, sin discutir».
\usepackage{float}

% --- Páginas apaisadas -----------------------------------------------
% pdflscape gira una página completa en el PDF. Imprescindible para la
% Tabla 3 (matriz de no conformidades) y la Tabla 4 (plan de acción).
\usepackage{pdflscape}

% --- Color (necesario para los enlaces de hyperref) ------------------
\usepackage{xcolor}
\definecolor{azuloscuro}{RGB}{0,51,102}

% =====================================================================
%  BIBLIOGRAFÍA — biblatex en estilo APA 7 con motor Biber
% =====================================================================
% csquotes es OBLIGATORIO antes de biblatex (gestiona las comillas).
\usepackage{csquotes}
\usepackage[backend=biber,style=apa,sortlocale=es_ES,natbib=true]{biblatex}
% Traduce al español los elementos fijos de las citas («y», «et al.», «s. f.»).
\DeclareLanguageMapping{spanish}{spanish-apa}
% Archivo con TODAS las referencias del informe.
\addbibresource{referencias.bib}

% =====================================================================
%  HYPERREF — SIEMPRE EL ÚLTIMO PAQUETE EN CARGARSE
% =====================================================================
% Convierte en enlaces clicables el índice, las referencias cruzadas y
% las URLs, y genera los marcadores (bookmarks) del PDF.
\usepackage[
  colorlinks=true,
  linkcolor=azuloscuro,     % enlaces internos (índice, \ref)
  citecolor=azuloscuro,     % citas bibliográficas
  urlcolor=azuloscuro,      % URLs externas
  bookmarks=true,
  bookmarksnumbered=true,
  breaklinks=true,
  % plainpages=false hace que las anclas del PDF usen el número de página tal
  % como se imprime (i, ii, 1, 2...). Junto con el \pagenumbering{roman} que
  % main.tex activa desde la portada, evita anclas duplicadas.
  plainpages=false,
  pdftitle={Auditoría socioambiental de la Mina Marlin — Informe EVI. 6},
  pdfsubject={Auditoría ambiental},
  pdfkeywords={auditoría socioambiental, Mina Marlin, Goldcorp, ISO 14001, ISO 19011, PHVA}
]{hyperref}
% Rompe las URLs largas de la bibliografía para que no se salgan del margen.
\usepackage{url}
\urlstyle{same}
